% ═══════════════════════════════════════════════════════════════
% LEHRERHINWEISE: Reaktion von Metallen mit Sauerstoff
% Thema 2.4 | Chemie Klasse 8 | Doppelstunde (90 Min)
% ═══════════════════════════════════════════════════════════════

\documentclass[11pt, a4paper]{article}

\usepackage[utf8]{inputenc}
\usepackage[T1]{fontenc}
\usepackage[ngerman]{babel}

\usepackage{helvet}
\renewcommand{\familydefault}{\sfdefault}

\usepackage{microtype}
\usepackage[margin=2cm, top=2cm, bottom=2cm]{geometry}
\usepackage{enumitem}
\usepackage{tabularx}
\usepackage{booktabs}
\usepackage{array}
\usepackage{multicol}

\usepackage{fancyhdr}
\usepackage{lastpage}

\usepackage{xcolor}
\usepackage{tcolorbox}
\tcbuselibrary{skins}

% ═══════════════════════════════════════════════════════════════
% FARBEN
% ═══════════════════════════════════════════════════════════════

\definecolor{chemie}{HTML}{2a7a4b}
\definecolor{hellgrau}{HTML}{F5F5F5}
\definecolor{mittelgrau}{HTML}{888888}
\definecolor{dunkelgrau}{HTML}{444444}
\definecolor{loesung}{HTML}{C62828}

\colorlet{fachfarbe}{chemie}

% ═══════════════════════════════════════════════════════════════
% KOPF- UND FUSSZEILE
% ═══════════════════════════════════════════════════════════════

\pagestyle{fancy}
\fancyhf{}
\renewcommand{\headrulewidth}{0.4pt}
\renewcommand{\footrulewidth}{0pt}
\lhead{\small Chemie Klasse 8 | Thema 2.4}
\rhead{\small\textcolor{fachfarbe}{\textbf{LEHRERHINWEISE}}}
\cfoot{\small\textcolor{mittelgrau}{Seite \thepage{} von \pageref{LastPage}}}

% ═══════════════════════════════════════════════════════════════
% BOXEN
% ═══════════════════════════════════════════════════════════════

\newtcolorbox{infobox}[1][]{
    colback=hellgrau,
    colframe=hellgrau,
    boxrule=0pt,
    arc=0pt,
    left=8pt, right=8pt, top=6pt, bottom=6pt,
    borderline west={3pt}{0pt}{fachfarbe},
    #1
}

\newtcolorbox{warnbox}{
    colback=white,
    colframe=loesung,
    boxrule=1pt,
    arc=0pt,
    left=8pt, right=8pt, top=6pt, bottom=6pt
}

% ═══════════════════════════════════════════════════════════════
% BEFEHLE
% ═══════════════════════════════════════════════════════════════

\newcommand{\abschnitt}[1]{%
    \vspace{10pt}%
    \noindent{\large\bfseries\textcolor{fachfarbe}{#1}}%
    \vspace{6pt}%
    \nopagebreak[4]%
}

\newcommand{\lsg}[1]{\textcolor{loesung}{\textbf{#1}}}

\setlength{\parindent}{0pt}
\setlength{\parskip}{4pt}

\newcolumntype{L}[1]{>{\raggedright\arraybackslash}p{#1}}
\newcolumntype{C}[1]{>{\centering\arraybackslash}p{#1}}

% ═══════════════════════════════════════════════════════════════
% DOKUMENT
% ═══════════════════════════════════════════════════════════════

\begin{document}

% ═══════════════════════════════════════════════════════════════
% HEADER
% ═══════════════════════════════════════════════════════════════

\begin{center}
    {\Large\bfseries\textcolor{fachfarbe}{Reaktion von Metallen mit Sauerstoff}}\\[0.3em]
    {\normalsize Lehrerhinweise | Chemie Klasse 8 | Thema 2.4 | DS (90 Min)}
\end{center}
\vspace{-0.3em}
\hrule
\vspace{0.8em}

% ═══════════════════════════════════════════════════════════════
% ÜBERBLICK
% ═══════════════════════════════════════════════════════════════

\abschnitt{Überblick}

\begin{infobox}
\textbf{Lernziele:}
\begin{itemize}[leftmargin=1.5em, itemsep=0pt]
    \item SuS definieren Oxidation als Reaktion mit Sauerstoff
    \item SuS formulieren Wort- und Formelgleichungen für Metalloxide
    \item SuS erklären das Gesetz der Massenerhaltung
    \item SuS vergleichen poröse (Mg) und dichte (Al) Oxidschichten
\end{itemize}
\end{infobox}

\textbf{Materialien:}
\begin{itemize}[leftmargin=1.5em, itemsep=0pt]
    \item AB-01-metalloxide (1× pro SuS)
    \item LB-01-metalloxide (optional als Referenz)
    \item Tablets/Laptops für LP und SIM
    \item Beamer für Video (Mg-Verbrennung)
\end{itemize}

\textbf{Digitale Ressourcen:}
\begin{itemize}[leftmargin=1.5em, itemsep=0pt]
    \item \texttt{LP-01-metalloxide.html} -- Interaktiver Lernpfad (8 Schritte)
    \item \texttt{SIM-01-oxidation.html} -- Simulation Metallverbrennung
    \item \texttt{MT-01-metalloxide.html} -- Stundenleistung (15 Min, 20 BE)
    \item Video: \texttt{vimeo.com/699344409} (Mg-Verbrennung)
\end{itemize}

% ═══════════════════════════════════════════════════════════════
% STUNDENVERLAUF
% ═══════════════════════════════════════════════════════════════

\abschnitt{Stundenverlauf (Doppelstunde)}

\begin{center}
\small
\begin{tabular}{|C{1cm}|L{2.8cm}|L{6cm}|L{3.5cm}|}
\hline
\textbf{Min} & \textbf{Phase} & \textbf{Inhalt} & \textbf{Material} \\
\hline
0--5 & Einstieg & Video Mg-Verbrennung; Leitfrage: Was passiert? & Beamer, Video \\
\hline
5--15 & Erarbeitung I & LP Schritte 1--2: Begriffe (Oxidation, Oxid, Element, Verbindung) & Tablet, LP \\
\hline
15--25 & Übung I & LP Schritt 3: Wortgleichungen; AB Aufg. 1--2 & AB, LP \\
\hline
25--40 & Erarbeitung II & LP Schritt 4: Formelgleichungen mit Wertigkeiten; AB Aufg. 3 & AB, LP \\
\hline
40--45 & \textit{Pause} & & \\
\hline
45--60 & Simulation & LP Schritt 5: SIM Metallverbrennung; AB Aufg. 4 (Flammenfarben) & Tablet, SIM \\
\hline
60--70 & Erarbeitung III & LP Schritte 6--7: Massenerhaltung, Korrosion; AB Aufg. 5--6 & AB, LP \\
\hline
70--80 & Sicherung & LP Schritt 8: Vergleich Mg/Al; AB Aufg. 7 & AB, LP \\
\hline
80--90 & Abschluss & Zusammenfassung; Merksätze ins Heft; Hausaufgabe & AB \\
\hline
\end{tabular}
\end{center}

% ═══════════════════════════════════════════════════════════════
% DIFFERENZIERUNG
% ═══════════════════════════════════════════════════════════════

\abschnitt{Differenzierung}

\begin{tabular}{|c|L{4cm}|L{8.5cm}|}
\hline
\textbf{Niveau} & \textbf{Zielgruppe} & \textbf{Anpassung} \\
\hline
$\star$ & Berufsreife (BR) & Fokus auf Wortgleichungen (Aufg. 1--2, 4); Formelgleichungen optional; mehr Zeit für Simulation \\
\hline
$\star\star$ & Mittlere Reife (MR) & Alle Aufgaben; Unterstützung bei Formelgleichungen durch Wertigkeitstabelle \\
\hline
$\star\star\star$ & Gymnasium (GY) & Zusatzaufgabe: Weitere Metalle (Ca, Zn); Transfer auf Redoxreaktionen \\
\hline
\end{tabular}

% ═══════════════════════════════════════════════════════════════
% HÄUFIGE FEHLER
% ═══════════════════════════════════════════════════════════════

\abschnitt{Häufige Fehler und Reaktionen}

\begin{warnbox}
\textbf{1. Verwechslung Oxidation/Oxid}\\
\textit{Fehler:} ``Oxidation ist ein Stoff''\\
\textit{Reaktion:} Oxidation = Vorgang (Reaktion), Oxid = Produkt (Stoff)

\vspace{6pt}

\textbf{2. Formelgleichung nicht ausgeglichen}\\
\textit{Fehler:} Mg + O$_2$ $\rightarrow$ MgO\\
\textit{Reaktion:} Atome zählen! Links 2 O, rechts 1 O $\rightarrow$ 2 Mg + O$_2$ $\rightarrow$ 2 MgO

\vspace{6pt}

\textbf{3. Massenerhaltung falsch verstanden}\\
\textit{Fehler:} ``Das Oxid ist schwerer, also ist die Masse nicht erhalten''\\
\textit{Reaktion:} Sauerstoff aus der Luft wird eingebaut! Gesamtmasse (Mg + O$_2$) = Masse (MgO)

\vspace{6pt}

\textbf{4. Passivierung nicht verstanden}\\
\textit{Fehler:} ``Aluminium oxidiert nicht''\\
\textit{Reaktion:} Al oxidiert sehr wohl, aber die Schicht ist dicht und schützt $\rightarrow$ Reaktion stoppt
\end{warnbox}

% ═══════════════════════════════════════════════════════════════
% MT-CODE
% ═══════════════════════════════════════════════════════════════

\abschnitt{Stundenleistung (MT)}

\begin{infobox}
\textbf{Differenziert nach Niveau:}\\[0.3em]
\begin{tabular}{lll}
$\star$ Basis & 6 Fragen & 15 BE \\
$\star\star$ Standard & +3 Fragen & 20 BE \\
$\star\star\star$ Erweitert & +2 Fragen & 25 BE
\end{tabular}
\end{infobox}

\textbf{Ablauf:}
\begin{enumerate}[leftmargin=1.5em, itemsep=0pt]
    \item SuS wählen \textbf{Platznummer} (P1--P30) und bestätigen
    \item SuS wählen \textbf{Niveau} ($\star$/$\star\star$/$\star\star\star$) und bestätigen
    \item Beide Wahlen sind nach Bestätigung \textbf{gesperrt}
    \item SuS bearbeiten Fragen ihres Niveaus (höhere Fragen informativ sichtbar)
    \item Abgabe $\rightarrow$ Ergebnis mit Notenpunkten $\rightarrow$ PDF speichern $\rightarrow$ Bildungscloud
\end{enumerate}

\textit{Hinweis:} Kein Unlock-Code nötig. Identifikation erfolgt über Platznummer.

% ═══════════════════════════════════════════════════════════════
% LÖSUNGEN
% ═══════════════════════════════════════════════════════════════

\newpage

\abschnitt{Lösungen zum Arbeitsblatt}

\textbf{Merksätze:}
\begin{itemize}[leftmargin=1.5em, itemsep=2pt]
    \item Merksatz 1: \lsg{Die Oxidation ist die chemische Reaktion eines Stoffes mit Sauerstoff. Bei der Oxidation entsteht ein Oxid und es wird Energie freigesetzt.}
    \item Merksatz 2: Ein \lsg{Oxid} ist eine chemische \lsg{Verbindung} aus einem Element und \lsg{Sauerstoff}.
\end{itemize}

\textbf{Aufgabe 1 (4 BE):}
\begin{tabular}{ll}
Oxidation & Reaktion mit Sauerstoff \\
Element & Stoff aus einer Atomsorte \\
Verbindung & Stoff aus versch. Atomen \\
Oxid & Verbindung mit Sauerstoff
\end{tabular}

\textbf{Aufgabe 2 (6 BE):}
\begin{enumerate}[label=\alph*), leftmargin=1.5em, itemsep=0pt]
    \item Mg + Sauerstoff $\rightarrow$ \lsg{Magnesiumoxid}
    \item Fe + Sauerstoff $\rightarrow$ \lsg{Eisenoxid}
    \item \lsg{Kupfer} + O$_2$ $\rightarrow$ Kupferoxid
    \item Al + Sauerstoff $\rightarrow$ \lsg{Aluminiumoxid}
    \item Zn + \lsg{Sauerstoff} $\rightarrow$ Zinkoxid
    \item \lsg{Calcium} + \lsg{Sauerstoff} $\rightarrow$ Calciumoxid
\end{enumerate}

\textbf{Aufgabe 3 (6 BE):}
\begin{enumerate}[label=\alph*), leftmargin=1.5em, itemsep=0pt]
    \item \lsg{2 Mg} + \lsg{O$_2$} $\rightarrow$ \lsg{2 MgO}
    \item \lsg{4 Al} + \lsg{3 O$_2$} $\rightarrow$ \lsg{2 Al$_2$O$_3$}
    \item \lsg{4 Fe} + \lsg{3 O$_2$} $\rightarrow$ \lsg{2 Fe$_2$O$_3$}
\end{enumerate}

\textbf{Aufgabe 4 (6 BE):}

a) Verbrennung von Magnesium:
\begin{tabular}{ll}
Vor Reaktion & \lsg{Silbrig glänzendes Magnesiumband} \\
Während & \lsg{Sehr helle, weiße Flamme; starke Hitze} \\
Nach Reaktion & \lsg{Weißes, sprödes Pulver (MgO)} \\
Energie & \lsg{Chemische Energie $\rightarrow$ Licht + Wärme (exotherm)}
\end{tabular}

b) Flammenfarben:
\begin{tabular}{lll}
Metall & Flamme & Oxid \\
Mg & \lsg{weiß (gleißend)} & \lsg{weiß} \\
Fe & \lsg{orange (Glut)} & \lsg{braun/rostrot} \\
Al & \lsg{silber-weiß} & \lsg{grau} \\
Cu & \lsg{grün-blau} & \lsg{schwarz}
\end{tabular}

\textbf{Aufgabe 5 (8 BE):}
\begin{enumerate}[label=\alph*), leftmargin=1.5em, itemsep=0pt]
    \item \lsg{Bei einer chemischen Reaktion bleibt die Gesamtmasse erhalten.}\\
    Gesetz: \lsg{Masse(Edukte) = Masse(Produkte)}
    \item Berechnung: \lsg{12\,g + 8\,g = \textbf{20\,g MgO}}
\end{enumerate}

\textbf{Aufgabe 6 (4 BE):}
\begin{enumerate}[label=\alph*), leftmargin=1.5em, itemsep=0pt]
    \item Beispiele: \lsg{Rost (Eisenoxid), Grünspan (Kupfercarbonat), angelaufenes Silber}
    \item Al rostet nicht: \lsg{Al bildet sofort eine dichte Oxidschicht, die das Metall schützt (Passivierung)}
\end{enumerate}

\textbf{Aufgabe 7 (6 BE):}

\begin{tabular}{lll}
Kriterium & Magnesium & Aluminium \\
Oxidschicht & \lsg{porös (durchlässig)} & \lsg{dicht (undurchlässig)} \\
Reaktion stoppt? & \lsg{Nein, vollständig} & \lsg{Ja, nur Oberfläche} \\
Fachbegriff & \lsg{Kettenreaktion} & \lsg{Passivierung}
\end{tabular}

Erklärung: \lsg{Bei Mg ist die Oxidschicht porös, Sauerstoff dringt weiter ein. Bei Al ist die Schicht dicht und blockiert weiteren Sauerstoff.}

\vfill
\hrule
\vspace{0.3em}
\textbf{Gesamtpunktzahl AB:} 37 BE

\end{document}
